\documentclass[10pt,a4paper]{amsart}
\usepackage[margin=19mm]{geometry}
\usepackage{amsmath,amssymb,mathtools}
\usepackage[scr=rsfs,cal=euler]{mathalfa}
\usepackage{xfrac}
\usepackage[unicode,hidelinks,bookmarks=false]{hyperref}
\newcommand{\ie}{i.e.\ }
\newcommand{\cf}{cf.\ }
\newcommand{\resp}{respectively\ }
\newcommand{\Subsection}{\subsection}
\providecommand{\nobreakdash}{\nobreak\hbox{-}}
\setlength{\emergencystretch}{2em}
\multlinegap0pt
\usepackage{bm}
\usepackage{bbm}
\usepackage{dsfont}
\usepackage{xspace}
\usepackage{cancel}
\usepackage{mathtools}
\usepackage{amsthm}
\makeatletter
\@ifundefined{proposition}{\newtheorem{proposition}{Proposition}}{}
\makeatother
\def\N{\mathbf{ N}}
\def\R{\mathbf{ R}}
\title[Enhanced CAD-Based Quantifier Elimination With Multiple Equa]{Enhanced CAD-Based Quantifier Elimination With Multiple Equational Constraints}
\author{James H. Davenport, Matthew England, Scott McCallum}
\date{Source: arXiv:2604.23873v2 (CC BY 4.0)}
\begin{document}
\maketitle

\noindent\emph{Source and licence.} Enhanced CAD-Based Quantifier Elimination With Multiple Equational Constraints. Version arXiv:2604.23873v2 (CC BY 4.0), distributed under the Creative Commons Attribution 4.0 International licence, as recorded in the arXiv licence metadata for this exact version and shown on the abstract page https://arxiv.org/abs/2604.23873v2. Full rights notice: licenses/arxiv-2604.23873-CC-BY-4.0.txt.\par\smallskip

\noindent\textit{Editorial scope.} Counted unit: the Proposition whose source label is \texttt{Invariance of solution vector cardinality} in the article (source0.tex lines 541--555, source label \texttt{Invariance of solution vector cardinality}), delivered together with the article's own complete original proof of it (source0.tex lines 557--637, one \texttt{proof} environment of 81 lines). No mathematics is written, repaired or re-derived: statement and proof are the article's own text line for line. Required context retained uncounted: none. Every definition, display and label of those lines is retained unchanged. Omitted from this package and not redistributed: 1--540, 556--556, 638--2255. Nothing inside a retained excerpt is omitted or altered: every retained body is delivered exactly as the article's lines are written, so no omission receipt of any kind accompanies this package. Citations: the retained excerpts carry no citation command at all, so no bibliography entry of the article is transcribed and no reference is redistributed. Reference adaptation: none was needed and none was made; every reference inside the delivered unit resolves inside it, so no unresolved reference or citation is produced. Printed slips kept literal: no printed word, symbol, hypothesis or number of the counted statement or of its proof was altered. Layout and class: the article's own class is \texttt{article} with packages including url, amsmath, amssymb, amsthm, pdfpages, verbatim, hyperref, combelow, ed, graphicx, xspace, physics, algorithm2e; the wrapper is the lane's pinned \texttt{amsart} editorial wrapper at 10pt on A4, which restates the theorem-like environments the retained text needs and adds the article's own macro definitions that the retained text needs (\texttt{\textbackslash N}, \texttt{\textbackslash R}), with extra wrapper packages (bm, bbm, dsfont, xspace, cancel, mathtools). The delivered numbering is the unit's own. Assets: no retained span contains a figure environment, a graphics-inclusion command, a PDF-page inclusion, a table or a TikZ picture; no image file is copied into this package, the package declares no asset dependency and no image file is redistributed with it. Licence: version arXiv:2604.23873v2 of this article is distributed under the Creative Commons Attribution 4.0 International licence, as recorded in the arXiv licence metadata for this exact version and shown on the abstract page https://arxiv.org/abs/2604.23873v2. A full rights notice is delivered at licenses/arxiv-2604.23873-CC-BY-4.0.txt. Characters: every retained excerpt is pure ASCII, so the delivered engine, XeLaTeX with its default Latin Modern text font, reports no missing character.
\begin{proposition} \label{Invariance of solution vector cardinality}
    Let $\mathcal{D}$ be a truth-invariant CAD of $\R^n$ 
    for $\phi$ and let $\mathcal{D}_s$ denote the CAD of $\R^s$ induced 
    by $\mathcal{D}$.
    Let $c$ be a cell of $\mathcal{D}_s$,
    and let $\alpha \in c$. Then the total number $\nu_{c, \alpha}$
    of distinct solution vectors $\beta \in \R^k$ for $\phi$
    over $\alpha$ is an element of 
    $\N \cup \{\infty\}$, constant as $\alpha$ varies in $c$:
    hence we may define $\nu_c$, the cardinality of the set of
    solution vectors associated with the points of $c$, to be this number.
    Moreover, in the case that $0 < \nu_c < \infty$, the $\nu_c$ distinct
    solution vectors for $c$ are expressible as continuous functions
    of the parameters in $c$.
\end{proposition}

\begin{proof}
    Recall that $1 \le k \le n-1$ and $s = n - k$ (from the preamble). 
    Define the proposition
    $\mathcal{P}(k)$ as follows: for every cell $c \in \mathcal{D}_{n-k}$
    and for every $\alpha \in c$, the number $\nu_{c, \alpha}$
    is a constant element of $\N \cup \{\infty\}$, as $\alpha$
    varies in $c$. We shall prove that $\mathcal{P}(k)$ is true for all
    $k$ in the range $1 \le k \le n-1$ by induction on $k$.

    The induction base is proved as follows.
    Let $c \in \mathcal{D}_{n-1}$. First suppose that the cylinder over $c$
    contains some sector-cell $\hat{c}$ of $\mathcal{D}$
    for which $\phi$ is true throughout $\hat{c}$.
    Take any $\alpha \in c$.
    Then, by the definition of sector-cell,
    there are infinitely many $\beta \in \R$ such that
    $(\alpha, \beta)$ lies in $\hat{c}$, hence for which
    $\phi(\alpha, \beta)$ is true.
    Therefore, $\nu_{c, \alpha} = \infty$.
    We have shown that $\nu_{c, \alpha}$ is constant (and infinite)
    for all $\alpha \in c$.
    Now suppose, on the other hand, that the cylinder over $c$
    contains no sector-cell for which $\phi$ is true.
    Take any $\alpha \in c$: then $\nu_{c, \alpha}$ equals the total number of section-cells
    of $\mathcal{D}$ over $c$ for which $\phi$ is true, a number clearly
    independent of $\alpha$. Again, we have shown that $\nu_{c, \alpha}$
    is constant for all $\alpha \in c$ (but in this case, it is finite).
    The induction base $\mathcal{P}(1)$ is now established.

    For the induction step, let $k > 1$ and assume $\mathcal{P}(k-1)$ is true
    (as induction hypothesis). We must show that $\mathcal{P}(k)$ holds.
    Let $c \in \mathcal{D}_{n-k}$.
    Suppose that the cylinder over $c$ contains some 
    cell $\hat{c} \in \mathcal{D}_{n-k+1}$ for which $\nu_{\hat{c}}$
    is infinite (noting that the notation $\nu_{\hat{c}}$
    is justified by the induction hypothesis).
    Then the number $\nu_{c, \alpha}$, for points $\alpha \in c$,
    is constant and infinite also.
    So, we may assume henceforth that the cylinder over $c$ contains
    no cell $\hat{c}$ in $\mathcal{D}_{n-k+1}$ for which $\nu_{\hat{c}}$
    is infinite. Two cases remain to be considered.
    Firstly, consider the case in which the cylinder over $c$
    contains some sector-cell $\hat{c} \in \mathcal{D}_{n-k+1}$
    for which $\nu_{\hat{c}}$ is positive
    (noting again that the notation is justified by the induction
    hypothesis). Then the number $\nu_{c, \alpha}$, for points
    $\alpha \in c$, is constant and infinite. The reason is that,
    with $\alpha \in c$ arbitrary,
    for each of the infinitely many $\beta_{s+1} \in \R$
    such that $(\alpha, \beta_{s+1}) \in \hat{c}$,
    there is a solution vector $(\beta_{s+2}, \ldots, \beta_n)$
    for $\phi$ associated with $(\alpha, \beta_{s+1})$.
    Secondly, consider the case in which the cylinder over $c$
    contains no such sector-cell. In this case, the number
    $\nu_{c, \alpha}$, for points $\alpha \in c$, equals the sum 
    over the sections $\hat{c}$ of $\mathcal{D}_{n-k+1}$
    of the finite quantities $\nu_{\hat{c}}$ (noting again the use of the
    induction hypothesis), a sum which is independent of $\alpha$.
    This completes the proof of the induction step,
    establishing the truth of $\mathcal{P}(k)$ for all $k$, with $1 \le k \le n-1$.

    Define the proposition $\mathcal{Q}(k)$ as follows:
    for every cell $c \in \mathcal{D}_{n-k}$,
    for which $0 < \nu_c < \infty$,
    the $\nu_c$ distinct solution vectors for $c$ are expressible
    as continuous functions of the parameters in $c$.
    Then $\mathcal{Q}(k)$ may be proved true for all $k$ in the range 
    $1 \le k \le n-1$ by induction on $k$,
    by analogy with (and suitable modification of)
    the proof for proposition $\mathcal{P}$ provided above.
    
    %For the induction base, let $c \in \mathcal{D}_{n-1}$.
    %By definition of CAD, the cells of $\mathcal{D}$ over $c$
    %consist of sections and sectors over $c$ (that is, a stack over $c$).
    %By hypothesis ($ \nu_c < \infty$), only sections are true for %$\phi$. 
    %By definition of section, the distinct solution vectors for $c$
    %are expressed as continuous functions of $x \in c$.
    %The induction step for $Q$ is proved by analogy with 
    %(and suitable modification of) the induction step for $P$
    %provided above.
\end{proof}

\end{document}
